% !TEX root = ../main.tex
% =====================================================================
% APPENDIX A -- THE INTERNSHIP ITSELF
%
% Host team, supervision, computing environment, the mission as it was
% given, the objectives with their success criteria, the constraints,
% how the plan evolved, and the skills acquired.
% Deliberately placed after the references so it does not interrupt the
% scientific story of Sections 1-7.
% =====================================================================

\section{Internship Context and Organisation}
\label{app:internship}

% ---------------------------------------------------------------------
\subsection{Host institution, team and supervision}
\label{app:host}

The internship took place at the IRCICA / CRIStAL laboratory, within the FOX research team, which
works on time series analysis across several application domains and on the hardware that runs it,
in particular microprocessors and microcontroller units. The subject was titled ``Time series
classification using machine learning for edge devices (microcontrollers)'', and the work divided
roughly into 70\,\% time series classification and 30\,\% edge-device considerations, which is the
balance reflected in this report.

I met my supervisor at least once a week, with day-to-day communication over Discord, and all code
was kept in a shared Git repository he had access to. I wrote the code alone, but every idea was
discussed before implementation: we read recent papers together, I proposed an approach and returned
with the measured results. That cycle produced the sequence of variants in \Cref{sec:methods}.

% ---------------------------------------------------------------------
\subsection{Computing environment}
\label{app:context:compute}

Code was written and small tests run on a local Windows machine with a 6\,GB RTX GPU. All training
reported here ran on the Grid'5000 testbed \citep{balouek2013grid5000}, using the Lille and Sophia
sites with the OAR job scheduler. The two environments are not identical -- only Grid'5000 produced
the numbers in this report -- and both are documented in full in the repository README
(\Cref{sec:setup:baselines}).

\begin{contribution}
I wrote the job submission scripts for both sites, and set up a monitoring system based on tmux and
a Telegram bot that sent the accuracy of each finished model to my phone, so a failed or stalled job
could be detected without keeping a session open.
\end{contribution}

% ---------------------------------------------------------------------
\subsection{The mission and how it was defined}
\label{app:mission}

\paragraph{Starting point.}
The subject began with a reading list from my supervisor: the \millet{} paper
\citep{early2024millet} and its public code, InterpGN \citep{wen2025interpgn}, VQShape
\citep{wen2024vqshape}, the MIL pooling literature \citep{ilse2018attention, li2021dsmil} and the TCN
paper \citep{bai2018tcn}. Reproducing \millet{} was step zero, not a contribution. One observation
from that phase shaped the whole methodology: re-running the published model did not reproduce the
published numbers exactly, which led to re-training every baseline under one identical configuration
(\Cref{sec:setup:baselines}) and, eventually, to the controlled comparison that closes
objective~O2.

\begin{priorwork}
The \millet{} code was reused without modification. The only change was to wrap it inside the
pipeline of \Cref{sec:setup:pipeline} so that their models and mine could be trained and evaluated
under identical conditions.
\end{priorwork}

\paragraph{Problem statement.}
The question the internship had to answer was whether a substantially smaller encoder can match the
accuracy of the \millet{} baseline while producing better explanations. This is not obvious in
either direction: reducing depth or width normally costs accuracy, and sharper explanations normally
cost accuracy as well, so improving on all three axes at once is not the default outcome.

% ---------------------------------------------------------------------
\subsection{Objectives}
\label{app:objectives}

\begin{table}[htbp]
\caption{The objectives set with my supervisor, their measurable success criteria, and the outcome.
\Cref{sec:conclusion} refers to this table rather than repeating it.}
\label{tab:objectives}
\begin{center}
\small
\begin{tabular}{p{0.20\linewidth} p{0.26\linewidth} p{0.44\linewidth}}
\toprule
\textbf{Objective} & \textbf{Success criterion} & \textbf{Outcome} \\
\midrule
O1. Reproduce the baselines
  & Compare architectures and results, and select one as the basis for interpretability
  & \textbf{Reached.} \millet{} was selected, and FCN, ResNet and InceptionTime were re-trained here
    under one identical configuration rather than quoted (\Cref{sec:setup:baselines}). \\
O2. Match the published \millet{} results
  & Re-run the \millet{} model and compare against the published numbers
  & \textbf{Reached.} Their released weights evaluate at 0.922; re-training the same architecture
    under the configuration of \Cref{tab:protocol} gives 0.887 on \webtraffic{}, and under their own
    published configuration 0.920 on \webtraffic{} and 0.8434 over \ucr{} against a published
    0.8445. The gap was the training budget, not the re-implementation (\Cref{sec:disc:reading}). \\
O3. Design a smaller encoder
  & $\geq$ 50\,\% fewer parameters at equal accuracy
  & \textbf{Reached, and exceeded.} 41{,}324 against 423{,}707 parameters at higher accuracy on
    \webtraffic{} -- a 90\,\% reduction (\Cref{tab:web_main}). \\
O4. Improve explanation quality
  & Higher AOPCR and NDCG@$n$ on \webtraffic{}
  & \textbf{Reached on \webtraffic{}.} Both metrics improve over the re-trained baseline, and the
    $\kappa$ ablation shows the gain comes from the intended mechanism (\Cref{sec:disc:topk}). Over
    the \ucr{} archive the proposed heads lose more datasets than they win
    (\Cref{sec:res:ucr}), so this is not a general claim. \\
O5. Build a reproducible benchmark
  & Any model reproducible from one command and one configuration file
  & \textbf{Reached.} Any encoder/pooling pair trains from a single YAML file, and every figure and
    table here is regenerated by one command. \\
O6. Test on a physical device
  & Measure performance of the final model on a microcontroller
  & \textbf{Not attempted.} The cost measurements indicate feasibility (41\,K parameters, 1.43\,MB,
    \Cref{tab:cost}), comparable in scale to the MLPerf Tiny reference model
    \citep{banbury2021mlperftiny}, but no on-device test was run. \\
\bottomrule
\end{tabular}
\end{center}
\end{table}

% ---------------------------------------------------------------------
\subsection{Constraints}
\label{app:constraints}

\begin{itemize}\itemsep2pt
  \item \textbf{Time.} Four months in total, covering the literature review, the implementation, the
        experiments and this report.
  \item \textbf{Compute.} Shared GPU nodes with a wall-clock limit per job, and a queue: a complete
        run over the 128 \ucr{} datasets takes about five hours for one model, and
        a job cannot always start when it is submitted. This is the direct reason most of the 72
        variants have a single seed (\Cref{sec:setup:protocol}).
  \item \textbf{Fairness of the comparison.} Every model had to be trained under one identical
        configuration, which meant that no model could be individually tuned.
  \item \textbf{Evaluation data.} NDCG@$n$ requires per-time-step labels, which only \webtraffic{}
        provides. Screening on \webtraffic{} was therefore not a choice but a constraint, and
        \Cref{sec:res:ucr} shows what it cost.
\end{itemize}

% ---------------------------------------------------------------------
\subsection{How the plan evolved}
\label{app:evolution}

The initial plan was to reproduce \millet{} and then shrink its encoder; two reading decisions
changed it. VQShape \citep{wen2024vqshape} explains predictions through learned shapelets, which we
judged more complex than the target required, so the MIL route was kept. Conversely, the
histopathology MIL literature proved directly relevant: the multi-scale idea of MS-DA-MIL
\citep{hashimoto2020msdamil} became the pyramid wrapper (\Cref{sec:var:pyramid}) and the two-stream
idea of DSMIL \citep{li2021dsmil} a pooling head (\Cref{sec:pool:dsmil}). Both transfers
underperformed, and the results explain why: a whole-slide image is a bag of thousands of patches
whose attention relies on a strong pre-trained feature extractor, whereas these series contain around
1{,}000 to 1{,}500 time steps and are encoded from scratch. That redirected the effort towards the
TCN-based encoders and the heads derived directly from limitations L1--L3. The final change was
methodological: after the \webtraffic{} results looked clean, the full \ucr{} experiments showed the
picture was incomplete (\Cref{sec:res:ucr}), which led to adding the multi-seed runs and the
controlled configuration comparison rather than reporting the first result alone.

% ---------------------------------------------------------------------
\subsection{Skills acquired}
\label{app:skills}

\paragraph{Scientific.}
Reproducing \millet{} required going beyond the equations in the paper into the authors' code, then
identifying why an honest re-implementation still missed the published number -- the training
configuration rather than the architecture. The same care shaped the protocol: re-training FCN,
ResNet and InceptionTime under one configuration is what allows any difference in \Cref{sec:results}
to be attributed to the architecture. I also learned to test a metric before trusting it, after
finding that the same model scores AOPCR 2.57 under one configuration and 13.27 under another.

\paragraph{Technical.}
Designing every encoder and head around one fixed tensor contract is the only reason 72 variants
could be trained without writing 72 model classes. Running them meant working on a shared HPC
platform, reserving GPU nodes on Grid'5000 with OAR across two sites so that a queue at one did not
stall the study. Making a single command rebuild every figure and table from the stored results
repeatedly caught numbers that had silently gone out of date.

\paragraph{Professional.}
The most difficult update to prepare was bringing my supervisor a discrepancy I had not yet solved
(objective~O2) instead of debugging privately until it matched; that turned a private problem into a
shared decision and produced a better experiment than I would have designed alone. The lesson I would
keep is the last one: before the 85-dataset \ucr{} experiments the \webtraffic{} numbers told a clean
story, and running the same models over 85 further datasets showed that the story does not fully
hold. Neither result is wrong; only the second is complete. That is the concrete form of not trusting
a result verified once, and it is why \Cref{sec:disc:limitations} is a substantive part of this
report.
